\documentclass{article}
\input{../../../../lib.sty}

\title{\texttt{fn make\_expr\_strptime}}
\author{\MichaelShoemateAuthor}
\date{}

\begin{document}

\maketitle

\contrib
Proves soundness of \rustdoc{transformations/make_stable_expr/namespace_str/expr_strptime/fn}{make\_expr\_strptime} in \asOfCommit{mod.rs}{f5bb719}.

\texttt{make\_expr\_strptime} returns a Transformation that parses strings in a dataframe column into temporal types.

\section{Hoare Triple}
\subsection*{Precondition}
\subsubsection*{Compiler-verified}
\begin{itemize}
    \item Argument \texttt{input\_domain} of type \texttt{WildExprDomain}
    \item Argument \texttt{input\_metric} of type \texttt{M}
    \item Generic \texttt{M} implements \texttt{OuterMetric}
    \begin{itemize}
        \item \texttt{OuterMetric} defines an associated type \texttt{InnerMetric} that must implement \texttt{UnboundedMetric}
    \end{itemize}
    \item \texttt{(ExprDomain, M)} implements \texttt{MetricSpace},
    \item \texttt{Expr} implements \texttt{StableExpr<M, M>}
\end{itemize}

\subsubsection*{Caller-verified}
None

\subsection*{Pseudocode}
\lstinputlisting[language=Python,firstline=2,escapechar=`]{./pseudocode/make_expr_strptime.py}

\subsection*{Postcondition}
\begin{theorem}
\validTransformation{\texttt{(input\_domain, input\_metric, expr, M)}}{\texttt{make\_expr\_strptime}}
\end{theorem}
Starting from line \ref{match-strategy},
\texttt{expr} is matched as if it were a strptime expression, or otherwise rejects the expression.

All preconditions for \texttt{make\_stable} on line \ref{t-prior} are compiler-verified,
therefore by the postcondition \texttt{t\_prior} is a valid transformation.

To prove that the output is a valid transformation,
we must first prove that the transformation on line \ref{strp-trans} is a valid transformation.

\begin{proof}
    \textbf{Data-Independent Errors}
    Line \ref{string-check} ensures that the input data is always a string type and
    line \ref{to-type-check} ensures that the output data is always a temporal type,
    which are necessary preconditions for the strptime function in Polars.
    However, even if these checks were ignored when they would fail,
    then all possible choices of input dataset would fail,
    resulting in data-independent errors.

    Line \ref{non-strict} disables raising an error for ordinary parse failures.
    Line \ref{ambiguous} replaces the ambiguity option with the constant \texttt{null}.
    Line \ref{to-type-check} rejects timezone-aware and nanosecond datetimes.
    In particular, the accepted datetime types do not localize wall-clock values
    into timezones, which could raise on nonexistent times during DST transitions.
    Line \ref{inferred-dtype} additionally rejects formats that would implicitly
    produce a timezone-aware output.

    Therefore, ordinary parse failures produce null values,
    without raising data-dependent errors.
\end{proof}

\begin{proof}
    \textbf{Appropriate Output Domain}
    Line \ref{out-series} constructs a series domain whose elements are of type \texttt{to\_type} or null.
    Each string in the input data may either parse into \texttt{to\_type} or fail and by \ref{non-strict} evaluate to null.

    Therefore the output domain is the same as the input domain,
    but where the active column domain's data type is nullable values of type \texttt{to\_type}.
\end{proof}

\begin{proof}
    \textbf{Stability Guarantee}
    strptime can be considered a 1-stable row-by-row transformation,
    as the date parsing is applied independently to each and every row.
    See \rustdoc{transformations/manipulation/fn}{make\_row\_by\_row} for a proof of the stability of row-by-row functions.
\end{proof}

Since it has been shown that both \texttt{t\_prior} and the strptime transformation are valid transformations,
then the preconditions for \texttt{make\_chain\_tt} are met (invoked via the right-shift operator shorthand).

\end{document}
